% figure
\definecolor{mycolor1}{rgb}{0.00000,0.30588,0.49020}%
\definecolor{mycolor2}{rgb}{0.1814, 0.4029 0.5784}%
\definecolor{mycolor3}{rgb}{0.3627, 0.5000, 0.6667}%
\definecolor{mycolor4}{rgb}{0.5441, 0.5971, 0.7549}%
\definecolor{mycolor5}{rgb}{0.7255, 0.6941, 0.8431}%
\definecolor{mycolor6}{rgb}{0.7357, 0.5639, 0.7294}%
\definecolor{mycolor7}{rgb}{0.7459, 0.4337, 0.6157}%
\definecolor{mycolor8}{rgb}{0.7561, 0.3035, 0.5020}%
\definecolor{mycolor9}{rgb}{0.7663, 0.1733, 0.3882}%
\definecolor{mycolor10}{rgb}{0.7765, 0.0431, 0.2745}%

\pgfplotscreateplotcyclelist{fycle}{%  %<------
	{mycolor1},%
	{mycolor2},%
	{mycolor3},%
	{mycolor4},%
	{mycolor5},%
	{mycolor6},%
	{mycolor7},%
	{mycolor8},%
	{mycolor9},%
	{mycolor10},%
}%

\begin{figure}[!htb] 
	\centering
	\tikzsetnextfilename{fig_close_loop_against_ki}
	\begin{tikzpicture}	
		\begin{axis}[%
			%name=ax1,
			width=0.85\linewidth,
			height=.3\linewidth,
			scale only axis,   
			xlabel={temps (\si{\second})},
			ylabel={Position (\si{\centi\meter})},
			ylabel near ticks,
			axis x line* = bottom,
			axis y line = left, 
			xmin=0,
			xmax=4,
			cycle list name = fycle,
			legend style={nodes={scale=0.75},at={(0.87,0.98)},anchor=north east},
			legend entries={0.01 SI, 0.08 SI, 9.1 SI},
			]
			%	
			\addlegendimage{mycolor1}
			\addlegendimage{mycolor3}
			\addlegendimage{mycolor10}
			\pgfplotsinvokeforeach{2,3,...,11}{
				\addplot+[
				smooth,
				]
				table [col sep=comma, x=t, y index=#1] {misc/chap_5/ki_influence.csv};
			}	
			%\addlegendentry{\SI{50}{\newton\per\meter}}					
		\end{axis}
		\begin{axis}[%
			axis y line={right,dashed},
			width=0.85\linewidth,
			height=.3\linewidth,
			scale only axis,  
			axis x line = none,
			ylabel={Force (\si{\newton})},
			xmin=0,
			xmax=4,
			]
			\addplot[
			color=Prune,
			dashed,
			smooth,
			]
			table [col sep=comma, x=t, y=u_in] {misc/chap_5/ki_influence.csv};
		\end{axis}
		%
	\end{tikzpicture}	
	\caption{Réponse du système en boucle fermée pour une entrée en force sinusoïdale (en tirets), en fonction du gain $K_i$ de l'action intégrale du contrôleur en admittance, variant de \SIrange{0.01}{9.1}{\meter\radian\per\newton\per\second\squared}.}
	\label{fig_close_loop_against_ki}
\end{figure}
%%
%